% =========================================================
% Proof: Product Rule
% Source: volume-iii/analysis/differentiation/notes/notes-algebra.tex
% Mode: standard
% =========================================================

\subsection*{Product Rule}
\label{prf:product-rule}

\begin{remark*}[Return]
\hyperref[thm:product-rule]{$\leftarrow$ Back to Theorem (Product Rule) in Notes}
\end{remark*}

\begin{proof}
Let $I \subseteq \mathbb{R}$ be an open interval, let $c \in I$, and let
$f, g : I \to \mathbb{R}$ be differentiable at $c$. Set $h := fg$. The claim
is that $h$ is differentiable at $c$ and $h'(c) = f'(c)g(c) + f(c)g'(c)$.

\ProofStep{1}{Form the difference quotient of $h$.}
For $x \in I$ with $x \neq c$:
\[
  \frac{h(x) - h(c)}{x - c}
  = \frac{f(x)g(x) - f(c)g(c)}{x - c}.
\]

\ProofStep{2}{Insert the auxiliary term $f(x)g(c)$.}
Add and subtract $f(x)g(c)$ in the numerator:
\[
  \frac{f(x)g(x) - f(x)g(c) + f(x)g(c) - f(c)g(c)}{x - c}.
\]
Factor each pair:
\[
  = f(x) \cdot \frac{g(x) - g(c)}{x - c}
  + g(c) \cdot \frac{f(x) - f(c)}{x - c}.
\]

\ProofStep{3}{Pass to the limit.}
By \ByThm{Sum Rule for Limits} and \ByThm{Product Rule for Limits}:
\[
  \lim_{x \to c} \frac{h(x) - h(c)}{x - c}
  = \left(\lim_{x \to c} f(x)\right)
    \left(\lim_{x \to c} \frac{g(x)-g(c)}{x-c}\right)
  + g(c) \left(\lim_{x \to c} \frac{f(x)-f(c)}{x-c}\right).
\]
Since $f$ is differentiable at $c$, it is continuous at $c$ by
\ByThm{Differentiability Implies Continuity}, so $\lim_{x \to c} f(x) = f(c)$.
The remaining limits equal $g'(c)$ and $f'(c)$ respectively by
\ByDef{Differentiability at a Point}. Therefore
\[
  h'(c) = f(c)g'(c) + g(c)f'(c).
\]
\end{proof}

\begin{remark*}[Proof shape]
The add-and-subtract trick inserts the term $f(x)g(c)$ to decouple the
increment of $f$ from the increment of $g$. This splits the difference
quotient of $fg$ into a sum of two terms, each containing exactly one
difference quotient. Continuity of $f$ at $c$ — supplied by differentiability
— handles the factor $\lim f(x) = f(c)$ in the first term.

\FlashProof{1. Form the difference quotient of $fg$.
  2. Insert $\pm f(x)g(c)$ to split into two difference quotients.
  3. Apply limit laws; use continuity of $f$ (from differentiability)
  for $\lim f(x) = f(c)$; differentiability of $f$ and $g$ gives $f'(c)$
  and $g'(c)$.}
\end{remark*}

\begin{remark*}[Dependencies]
\begin{itemize}
  \item \hyperref[def:derivative-at-point]{Definition of Differentiability at a Point}:
    Identifies the difference quotient limits of $f$ and $g$ as $f'(c)$
    and $g'(c)$.
  \item \hyperref[thm:diff-implies-cont]{Differentiability Implies Continuity}:
    Supplies $\lim_{x \to c} f(x) = f(c)$ in Step 3.
  \item \hyperref[thm:limit-sum]{Sum Rule for Limits}:
    Distributes the limit across the two-term sum.
  \item \hyperref[thm:limit-product]{Product Rule for Limits}:
    Separates the limit of $f(x) \cdot \frac{g(x)-g(c)}{x-c}$ into a
    product of limits.
\end{itemize}

\FlashUses{def:derivative-at-point, thm:diff-implies-cont, thm:limit-sum, thm:limit-product}
\end{remark*}